\clearpage\fancyhead[L]{\small KOMET.}\hypertarget{native-340-lower}{}\label{native-340-lower}
\begin{center}\textbf{KOMET.}\par\vspace{4pt}\hrule\vspace{4pt}\fontsize{10}{12}\selectfont Monatliche Correspondenz. Band XXIV. S. 304\,..\,307. 1811 September.\par\vspace{4pt}\hrule\vspace{4pt}\end{center}

\par\indent ---\;---\;--- Professor \textsc{Gauss} hat uns interessante Ergebnisse seiner Beobachtungen und Rechnungen mitgeteilt. Aus einigen seiner Briefe geben wir die einschlägigen Stellen wieder. ---\;---\;--- Göttingen, am 8. September. „Den Kometen habe ich zuerst am 22. August gesehen. Da die Stadt den Blick von der Sternwarte nach Norden stark einschränkt, konnte ich ihn erst am 3. September zum ersten Mal beobachten. Ich habe versucht, seine Position aus Abständen zu bestimmen, die ich mit dem Sextanten gemessen habe. Mit dem schönen Stativ, das Sie uns über den geschickten Instrumentenbauer \textsc{Körner} in Weimar besorgt haben, ist das eine sehr bequeme und, wie es scheint, auch verhältnismäßig sehr genaue Beobachtungsmethode. Die Beobachtungen vom 3. September habe ich noch nicht reduziert. Die Beobachtungen der folgenden Tage, bei denen zum Vergleich immer $\alpha$ Aurigae und $\alpha$ Lyrae herangezogen wurden, habe ich mit aller Sorgfalt berechnet und folgende Ergebnisse gefunden.
\Needspace{9\baselineskip}
\begin{center}\fontsize{10}{12}\selectfont\setlength{\tabcolsep}{4pt}\renewcommand{\arraystretch}{1.15}\begin{tabular}{r|c|c|c}
1811 & \shortstack{Mittlere Zeit\\in Göttingen} & \shortstack{Rektaszension\\des Kometen} & \shortstack{Nördliche\\Deklination}\\\hline
Sept. 4 & $8^{\mathrm u}\,28^{\mathrm m}\,47^{\mathrm s}$ & $158^\circ\,25'\,24''$ & $39^\circ\,18'\,2''$\\
6 & $8\;48\;38$ & $160\;23\;16$ & $40\;14\;16$\\
7 & $8\;57\;6$ & $161\;26\;10$ & $40\;41\;54$\\
\end{tabular}\end{center}

\SourceLeaf{341}\par\indent Nach diesen und neun Beobachtungen von Herrn von \textsc{Zach} habe ich meine parabolischen Elemente\footnote{Monatl. Corresp., Augustheft 1811, S. 180. [S. 326 d. B.]} verbessert und dabei die folgenden Werte erhalten.
\Needspace{10\baselineskip}
\begin{center}\fontsize{10}{12}\selectfont\setlength{\tabcolsep}{4pt}\renewcommand{\arraystretch}{1.15}\begin{tabular}{l}
Zeit des Periheldurchgangs 1811 Sept. 12\\
\quad $5^{\mathrm u}\,21^{\mathrm m}\,15^{\mathrm s}$ Mittlere Zeit in Göttingen\\
\end{tabular}\end{center}
\begin{center}\fontsize{10}{12}\selectfont\setlength{\tabcolsep}{4pt}\renewcommand{\arraystretch}{1.15}\begin{tabular}{llr}
Logarithmus des kleinsten Abstands & .\;.\;.\;.\;.\;.\;.\;. & $0.017060$\\
Länge des Perihels & .\;.\;.\;.\;.\;.\;.\;.\;.\;.\;.\;.\;. & $75^\circ\,17'\,34''$\\
Länge des aufsteigenden Knotens & .\;.\;.\;.\;.\;.\;.\;.\;. & $140\;24\;13$\\
Neigung der Bahn & .\;.\;.\;.\;.\;.\;.\;.\;.\;.\;.\;.\;.\;. & $73\;7\;16$\\
\end{tabular}\end{center}

\par\indent Es scheint bisher noch keinen sicheren Hinweis darauf zu geben, dass die Bahn elliptisch ist. Ich habe mir bisher nur ein paar Mal die Zeit genommen, den Kometen mit einem großen Teleskop zu betrachten. Am 7. September war die Form des Schweifs sehr bemerkenswert. Er bog sich in zwei Ästen vom Kometen weg. Diese gingen aber nicht vom Kometen selbst aus, sondern hingen in geringer Entfernung von ihm zusammen. Ein dunkler Zwischenraum trennte sie vom Kometen, sodass sie ihn wie eine Parabel ihren Brennpunkt umschlossen. Gewiss ist dieser Komet in mehr als einer Hinsicht einer der bemerkenswertesten, die je beobachtet wurden. Ich konnte ebenso wenig wie Professor \textsc{Harding} eine Spur eines eigentlichen Kerns im Kometen wahrnehmen.“ ---\quad Göttingen, am 16. September. Da Sie den Kometen jetzt im Meridian beobachten, werde ich meine Beobachtungen heute Abend abschließen. Hier ist, was ich seit meinem letzten Brief erhalten habe.
\Needspace{10\baselineskip}
\begin{center}\fontsize{10}{12}\selectfont\setlength{\tabcolsep}{4pt}\renewcommand{\arraystretch}{1.15}\begin{tabular}{r|c|c|c}
 & \shortstack{Mittlere Zeit\\in Göttingen} & \shortstack{Rektaszension\\des Kometen} & \shortstack{Nördliche\\Deklination}\\\hline
Sept. 9 & $8^{\mathrm u}\,37^{\mathrm m}\,53^{\mathrm s}$ & $163^\circ\,35'\,5''$ & $41^\circ\,39'\,13''$\\
14 & $8\;45\;5$ & $169\;44\;58$ & .\quad.\quad.\quad.\\
15 & $7\;39\;57$ & $171\;3\;58$ & $44\;24\;57$\\
16 & $7\;57\;16$ & $172\;31\;4$ & $44\;51\;23$\\
\end{tabular}\end{center}

\Needspace{24\baselineskip}\par\indent Für den Lauf des Kometen bis Ende Dezember habe ich aus meinen verbesserten Elementen die folgende Ephemeride berechnet.
\begin{center}\fontsize{10}{12}\selectfont\setlength{\tabcolsep}{3pt}\renewcommand{\arraystretch}{1.15}\begin{tabular}{rr|c|c|c|c}
\multicolumn{2}{c|}{\shortstack{Mittlere Zeit\\in Göttingen}} & \shortstack{Rektaszension} & \shortstack{Nördliche\\Deklination} & \shortstack{Log. des Abstands\\des Kometen\\von der Erde} & \shortstack{Helligkeit}\\\hline
1811 Sept. 12 & $5^{\mathrm u}$ & $166^\circ\,54'$ & $42^\circ\,58'$ & $0.2079$ & $0.355$\\
17 & $0$ & $173\;46$ & $45\;15$ & $0.1836$ & $0.394$\\
22 & $6$ & $182\;2$ & $47\;17$ & $0.1594$ & $0.432$\\
.27 & $6$ & $192\;19$ & $48\;48$ & $0.1366$ & $0.465$\\
Okt. 2 & $7$ & $203\;24$ & $49\;34$ & $0.1160$ & $0.490$\\
7 & $7$ & $215\;58$ & $49\;7$ & $0.1004$ & $0.501$\\
12 & $8$ & $228\;45$ & $47\;16$ & $0.0911$ & $0.492$\\
22 & $9$ & $251\;23$ & $39\;46$ & $0.0963$ & $0.421$\\
Nov. 1 & $10$ & $267\;58$ & $30\;32$ & $0.1312$ & $0.310$\\
11 & $11$ & $279\;38$ & $22\;18$ & $0.1851$ & $0.208$\\
21 & $12$ & $288\;13$ & $14\;55$ & $0.2415$ & $0.138$\\
Dez. 1 & $13$ & $294\;54$ & $9\;42$ & $0.2976$ & $0.093$\\
11 & $14$ & $300\;24$ & $5\;49$ & $0.3487$ & $0.064$\\
21 & $15$ & $305\;7$ & $2\;55$ & $0.3938$ & $0.046$\\
31 & $16$ & $309\;16$ & $0\;48$ & $0.4330$ & $0.034$\\
\end{tabular}\end{center}

\Needspace{16\baselineskip}\par\indent Ein ganz flüchtiger Vergleich meiner Beobachtungen und der ersten Beobachtungen von \textsc{Olbers} mit meinen Elementen hat mir folgende Ergebnisse geliefert;
\SourceLeaf{342}
\begin{center}\fontsize{10}{12}\selectfont\setlength{\tabcolsep}{7pt}\renewcommand{\arraystretch}{1.15}\begin{tabular}{lr|r|r}
 & & AR. & Dekl.\\\hline
Aug. & 23 & $+127''$ & $+18''$\\
Sept. & 4 & $+30$ & $-9$\\
 & 6 & $+9$ & $-11$\\
 & 7 & $-54$ & $+20$\\
 & 9 & $-65$ & $-64$\\
 & 14 & $-140$ & .\;.\;.\;.\\
 & 15 & $-195$ & $-76$\\
\end{tabular}\end{center}

\par\indent Ist der hier erkennbare, schnell wachsende Unterschied in der Rektaszension begründet, so wäre ich geneigt, darin bereits ein Anzeichen für die Elliptizität der Bahn zu sehen.“
